\documentclass[11pt]{amsart}

\usepackage{amsmath,amssymb,amsthm,mathtools}
\usepackage{geometry}
\usepackage{booktabs}
\usepackage{hyperref}
\usepackage{xcolor}
\usepackage{enumitem}
\geometry{margin=1in}
\hypersetup{colorlinks=true,linkcolor=blue!70!black,citecolor=blue!70!black,
  urlcolor=blue!70!black,
  pdftitle={A Faithful Krein-Space Realization of the Weil Quadratic Form},
  pdfauthor={David Alejandro Trejo Pizzo}}

\newtheorem{theorem}{Theorem}[section]
\newtheorem{proposition}[theorem]{Proposition}
\newtheorem{lemma}[theorem]{Lemma}
\newtheorem{corollary}[theorem]{Corollary}
\theoremstyle{definition}
\newtheorem{definition}[theorem]{Definition}
\newtheorem{hypothesis}[theorem]{Hypothesis}
\theoremstyle{remark}
\newtheorem{remark}[theorem]{Remark}

\newcommand{\re}{\operatorname{Re}}
\newcommand{\tr}{\operatorname{tr}}
\newcommand{\spec}{\operatorname{spec}}
\newcommand{\half}{\tfrac12}
\newcommand{\PW}{\mathit{PW}_d}
\newcommand{\tform}{\mathfrak t}
\newcommand{\tplus}{\mathfrak t_+}
\newcommand{\tminus}{\mathfrak t_-}
\newcommand{\ttilde}{\widetilde{\mathfrak t}_-}
\newcommand{\Hplus}{H_+}
\newcommand{\rmar}{\mathfrak a}
\newcommand{\rmq}{\mathfrak q}
\newcommand{\rmp}{\mathfrak p}
\newcommand{\Top}{\mathcal T}
\newcommand{\Kop}{K}
\newcommand{\Norm}[1]{\lVert #1 \rVert}

\title[A Kre\u{\i}n--Space Weil Realization]{A Faithful Kre\u{\i}n--Space Realization of the Weil Quadratic Form, with a No-Go Theorem for the Zero-Side Attack on Positivity}
\author{David Alejandro Trejo Pizzo}
\thanks{Independent Researcher. \texttt{dtrejopizzo@gmail.com}}


\begin{document}
\begin{abstract}
Weil's criterion identifies the Riemann Hypothesis (RH) with the positive
semidefiniteness of the Weil functional $\tform$ read from the Riemann--Weil explicit
formula. We give a rigorous operator-theoretic realization of $\tform$ on the
classical band-limited (Paley--Wiener) test class $\PW$, $d>\half$, in three layers.
\textbf{(A) Semibounded realization.} Modulo a single, precisely isolated hypothesis
\textnormal{(H)}---a uniform short-interval zero-density statement strictly weaker than
RH---the form $\tform$ is closable and defines a self-adjoint operator $\Top$ that is
bounded below, with $\inf\spec(\Top)$ a finite, convergent quantity; consequently
$\mathrm{RH}\iff \inf\spec(\Top)\ge0$. Four supporting lemmas are unconditional.
\textbf{(B) Kre\u{\i}n structure.} We prove the exact identity $\tform=E^{*}JE$, where
$E$ is evaluation at the (near-line) zeros and $J$ is a fundamental symmetry whose
negative directions are exactly the off-line zeros; the construction is therefore a
Kre\u{\i}n-space problem (indefinite metric), not a de~Branges (a priori positive)
one. The angular operator $\Kop$ of the sampling range yields a single-axis dictionary:
$\Norm{\Kop}<\infty\iff$ finiteness of the bottom, and $\Norm{\Kop}\le1\iff\mathrm{RH}$.
\textbf{(C) No-go theorem.} We prove that the explicit dependence of $\Kop$ on the
off-line depths carries no leverage toward the sign: the zero-side expression and the
arithmetic side of the explicit formula are the same functional, so bounding
$\Norm{\Kop}\le1$ from the zero side is identically the problem of proving Weil
positivity. We situate the construction relative to the Connes--Consani archimedean
trace approach and the Suzuki de~Branges completion: ours is strongly analogous,
the distinguishing feature being its unconditional, explicitly Kre\u{\i}n, finite
sampling ladder. We do \emph{not} claim progress past the sign barrier; we claim a
precise localization of it. This paper closes the explicit-formula/zero-side line of
this line of work and isolates the only remaining route to RH (structural positivity).
\end{abstract}

\maketitle



\tableofcontents
\newpage

% ════════════════════════════════════════════════════════════════════════
\section{Introduction}\label{sec:intro}
\subsection{Weil's criterion and the sign problem}
Let $\xi(s)=\half s(s-1)\pi^{-s/2}\Gamma(s/2)\zeta(s)$ be the completed Riemann zeta
function, with nontrivial zeros $\rho$. For a suitable even test function the
Riemann--Weil explicit formula equates a sum over zeros with an archimedean term, a
pole term, and a sum over primes. Polarizing, one obtains a Hermitian quadratic
functional $\tform(F)$ with the property
\begin{equation}\label{eq:weil}
\boxed{\ \mathrm{RH}\iff \tform(F)\ge0 \ \text{ for all admissible } F.\ }
\end{equation}
This is Weil's positivity criterion \cite{Weil1952,Bombieri}. In this formulation the
difficulty of RH is \emph{entirely the sign} of $\tform$: the right-hand side of the
explicit formula is indefinite. On-line zeros contribute $|F(\gamma)|^{2}\ge0$, while
off-line zeros (should they exist) contribute terms of \emph{either} sign. Proving the
sign is proving RH; no weaker statement about the \emph{magnitude} of $\tform$ suffices.
This paper makes that precise and operator-theoretic, and then proves a structural
theorem explaining why the most natural (zero-side) realization cannot, by itself,
decide the sign.

\subsection{The three layers and the main results}
We work on the band-limited class $\PW$ (Paley--Wiener, exponential type $d>\half$,
$L^{2}$ on the line). Band-limited functions are entire, so the off-axis evaluations
$F(\gamma+ib)$ forced by off-line zeros are automatic and $L^{2}$-controlled; this is
the technically clean home for Weil's criterion.

\begin{theorem}[Semibounded realization; \S\ref{sec:thmA}]\label{thm:A}
Assume \textnormal{(H)} \textnormal{(Definition~\ref{def:H})}. On $\PW\cap H^{1}$,
$d>\half$, the Weil form $\tform$ is closable in $\Hplus$ and its self-adjoint operator
$\Top$ satisfies $\inf\spec(\Top)\ge 1-4C>-\infty$ for a finite $C=C(d,\delta,\eta)$.
Consequently $\mathrm{RH}\iff\inf\spec(\Top)\ge0$.
\end{theorem}

\begin{theorem}[Kre\u{\i}n identity and the single axis; \S\S\ref{sec:krein}--\ref{sec:angular}]\label{thm:B}
Unconditionally, $\tform=E^{*}JE$ with $J=J^{*}$, $J^{2}=I$. The angular operator
$\Kop$ of the closed sampling range is defined (type-density uniqueness), and
\[
\Norm{\Kop}<\infty\iff \text{(finite bottom, B-2)},\qquad
\Norm{\Kop}\le1\iff \mathrm{RH},
\]
with $\Norm{\Kop}\ge e^{d/2}>1$ from the magnitude side.
\end{theorem}

\begin{theorem}[No-go; \S\ref{sec:nogo}]\label{thm:C}
The off-line depths $\{b_\rho\}$ enter $\Norm{\Kop}$ only through $\tform$ itself. Hence
bounding $\Norm{\Kop}\le1$ from the zero side is identically Weil positivity, and every
zero-position-free, sign-agnostic bound gives $\Norm{\Kop}^{2}\gtrsim e^{d}$ (the wrong
side of $1$). The zero-side realization, being defined in terms of the unknown off-line
zeros, cannot attack the sign.
\end{theorem}

\subsection{Relation to prior work and honest scope}
The reduction ``$\Norm{\Kop}\le1\iff$ RH'' is strongly analogous to the device of
Connes--Consani \cite{ConnesConsani}, who root Weil positivity in the compression of the
scaling action onto Sonin's space and control a residual operator by an
eigenvalue-below-$1$ argument. The completion of the Weil form into a reproducing-kernel
space is the theme of Suzuki \cite{Suzuki}, who obtains a \emph{de~Branges} space under
RH. Our construction differs in being \emph{unconditional} (no RH assumed; the metric is
genuinely indefinite, hence Kre\u{\i}n, not de~Branges) and in exposing an explicit
finite sampling ladder. We do \emph{not} claim an identification of our $\Kop$ with the
Connes--Consani operator; Theorem~\ref{thm:C} in fact explains why the zero-side
realization, unlike the archimedean one, cannot serve as a computational handle.

Theorem~\ref{thm:A} is conditional on (H); (H) is much weaker than RH but is not
established unconditionally by current methods (\S\ref{sec:Hstatus}). Theorems~\ref{thm:B}
and \ref{thm:C} are unconditional. \textbf{None of the three proves RH}, and
Theorem~\ref{thm:C} proves that this particular route cannot. The value of the paper is a
clean, faithful reformulation together with a rigorous map of the wall, and the
identification of the only remaining route (structural positivity, \S\ref{sec:open}).

% ════════════════════════════════════════════════════════════════════════
\section{Setup and notation}\label{sec:setup}

\paragraph{Test class.} For $d>\half$ let
$\PW=\{F\in L^{2}(\mathbb R): \widehat F\in L^{2}[-d,d]\}$ be the Paley--Wiener space
(entire functions of exponential type $\le d$ in $L^{2}$), with
$\widehat F(\xi)=\int F(t)e^{-it\xi}\,dt$. We use the dense domain
$\mathcal D:=\PW\cap H^{1}$ (so $\widehat F\in H^{1}[-d,d]\hookrightarrow C^{0}$). Two
classical inequalities recur:
\begin{equation}\label{eq:bernstein}
\text{(Bernstein)}\quad \Norm{F'}_{2}\le d\,\Norm{F}_{2},\qquad
\text{(Nikolskii)}\quad \Norm{F}_{\infty}\le C_{d}\Norm{F}_{2}.
\end{equation}

\paragraph{Zeros, depths, partners.} Write a nontrivial zero as
$\rho=\half+b_\rho+it_\rho$, $b_\rho,t_\rho\in\mathbb R$, so $b_\rho=0\iff\rho$ on the
critical line. The functional equation gives the symmetric quartet
$\{\rho,\bar\rho,1-\rho,1-\bar\rho\}$. Define the reflected partner involution
\begin{equation}\label{eq:sigma}
\sigma(\rho):=1-\bar\rho=\half-b_\rho+it_\rho,\qquad
\sigma^{2}=\mathrm{id},\quad b_{\sigma(\rho)}=-b_\rho,\ t_{\sigma(\rho)}=t_\rho.
\end{equation}
The zero density is $\rho(T)=\tfrac1{2\pi}\log\tfrac T{2\pi}$, and $S(T)=N(T)-\bar
N(T)=O(\log T)$ unconditionally. The Nyquist cell has length $1/d$ and contains
$\sim\rho(T)/d\to\infty$ zeros (over-sampling).

\paragraph{Weil functional.} For $F\in\mathcal D$ let $h(z):=F(z)\overline{F(\bar z)}$.
The explicit formula gives
\begin{equation}\label{eq:EF}
\tform(F):=\sum_\rho h(\text{spectral point of }\rho)
=\underbrace{\rmar(F)}_{\text{arch.}}+\underbrace{\rmq(F)}_{\text{pole}}
-\underbrace{\rmp(F)}_{\text{primes}},
\end{equation}
the spectral point of $\rho=\half+b_\rho+it_\rho$ being $t_\rho+ib_\rho$ (on-line zeros
sampled on the real axis). Grouping by partners \eqref{eq:sigma} gives the on/off
decomposition
\begin{equation}\label{eq:onoff}
\tform(F)=\underbrace{\sum_{b_\rho=0}|F(t_\rho)|^{2}}_{\tplus(F)\,\ge0}
\;+\;\underbrace{\sum_{\text{off-line pairs}}2\re\!\big[F(t_\rho+ib_\rho)\overline{F(t_\rho-ib_\rho)}\big]}_{\tminus(F)\ \text{indefinite}} .
\end{equation}
With $\ttilde(F):=\sum_{b_\rho\ne0}|F(t_\rho+ib_\rho)|^{2}$ one has
$|\tminus(F)|\le 2\,\ttilde(F)$ (sharpened to an exact projection identity in
\S\ref{sec:angular}).

\paragraph{Realization space.} Let $\Hplus$ be the completion of $\mathcal D$ in the
seminorm $\Norm{F}_{\Hplus}^{2}:=\tplus(F)$ (the on-line sampling space). $\Top$ is the
self-adjoint operator associated with the form $\tform$ on $\Hplus$.

% ════════════════════════════════════════════════════════════════════════
\section{The four unconditional lemmas}\label{sec:lemmas}

\begin{lemma}[Integration by parts]\label{lem:A}
For $F\in\mathcal D$ and $g:=|F|^{2}$,
$E_g:=\sum_{b_\rho=0} g(t_\rho)-\int g\,\rho\,dt=-\int g'(t)\,S(t)\,dt$.
\end{lemma}
\begin{proof}
$\widehat F\in H^{1}\hookrightarrow C^{0}$, and IBP in $\xi$ gives $|F(t)|=O(1/t)$, so
$g=O(t^{-2})$ and $g\log t\to0$; the boundary terms $[gS]_{\pm\infty}$ vanish, and
$dS=dN-\rho\,dt$ has locally bounded variation, so Stieltjes IBP yields the claim.
\end{proof}

\begin{lemma}[The quadrature error is $O(\tplus)$]\label{lem:B}
$|E_g|\le 2d\,(\sup_{\mathrm{supp}\,g}|S|)\,\Norm{F}_2^{2}=O(\rho\Norm{F}_2^{2})=O(\tplus(F))$.
\end{lemma}
\begin{proof}
$g'=2\re(F'\overline F)$, so $\int|g'|\le 2\Norm{F'}_2\Norm{F}_2\le 2d\Norm{F}_2^2$ by
Bernstein \eqref{eq:bernstein}. With $|S|=O(\log T_0)$ on $\mathrm{supp}\,g$ and the tail
controlled by $g'=O(t^{-3})$, $|E_g|\le C_d(\log T_0)\Norm{F}_2^2$. As the on-line main
term is $\sim\rho\Norm{F}_2^2\sim\tfrac{\log T_0}{2\pi}\Norm{F}_2^2$, this is $O(\tplus)$.
\end{proof}

\begin{remark}\label{rem:Onoto}
B-2 (a finite bottom) needs only $E_g=O(\tplus)$; the sharp $o(\tplus)$---which would
require Montgomery pair-correlation / RH-level input---is \emph{not} needed. This is the
single most important calibration of this line of work: we need an $O$, not an $o$.
\end{remark}

\begin{lemma}[Vertical $L^{2}$ bound]\label{lem:C}
For $F\in\PW$ and $|b|<\half$,
$\Norm{F(\cdot+ib)}_2^{2}=\int_{-d}^{d}|\widehat F(\xi)|^{2}e^{-2b\xi}d\xi\le e^{2d|b|}\Norm{F}_2^{2}\le e^{d}\Norm{F}_2^{2}.$
\end{lemma}
\begin{proof}
Parseval on the shifted line, then $e^{-2b\xi}\le e^{2d|b|}\le e^{d}$ on $|\xi|\le d$,
$|b|<\half<d$. Exact, no RH.
\end{proof}
The amplification $\kappa(b):=e^{2d|b|}$ ($>1$ for $b\ne0$) is the quantitative source of
the lower bound $\Norm{\Kop}\ge e^{d/2}$ in \S\ref{sec:angular}.

\begin{lemma}[Off-line sum bound; Plancherel--P\'olya]\label{lem:D}
$\ttilde(F)=\sum_{b_\rho\ne0}|F(t_\rho+ib_\rho)|^{2}\le C_d\,\rho\,\Norm{F}_2^{2}.$
\end{lemma}
\begin{proof}
Off-line spectral points lie in the strip interior $|b_\rho|<\half<d$ (margin
$d-\half>0$). The band-limited Plancherel--P\'olya/Bessel inequality, valid for
arbitrary (possibly clustered) sequences in a strip interior, gives
$\sum_n|F(w_n)|^{2}\le B\Norm{F}_2^{2}$ with $B\asymp_d\sup_x\#\{n:\re w_n\in[x,x+\tfrac1d]\}$.
The pointwise bound $|F(w)|^{2}\le C_d\int_{|x-\re w|<1/d}|F|^{2}$ summed over the cell
count gives $B=O(\rho)$ by Riemann--von~Mangoldt and $S=O(\log)$.
\end{proof}

\begin{remark}[Why $O(\rho)$ and not $o(\rho)$]
Zero-density makes off-line zeros globally sparse, but they may cluster locally; the cell
count is bounded by the total $O(\rho)$, not by the off-line density. Replacing it by the
latter would give $o(\tplus)=$ RH---forbidden by possible clustering.
\end{remark}

% ════════════════════════════════════════════════════════════════════════
\section{Coercivity, the depth split, and the hypothesis (H)}\label{sec:H}

By Lemmas~\ref{lem:C}--\ref{lem:D}, $\ttilde\le C_d\rho\Norm{F}_2^2$. To conclude
$\ttilde\le C\tplus$ (which is B-2) we need \emph{coercivity}
$\tplus(F)\ge c\,\rho\,\Norm{F}_2^{2}$, a lower sampling bound. This is not free: by the
Beurling gap theorem the on-line zeros are a set of sampling for $\PW$ only if they have
no gap exceeding $\pi/d$; the density condition $D^{-}\ge d/\pi$ is necessary but not
sufficient. A single Nyquist gap containing a bump of $F$ destroys coercivity.

\paragraph{Depth split.} Coercivity is needed only where $\ttilde>0$, i.e. near off-line
zeros. Split off-line zeros by depth at a threshold $\delta$: \emph{shallow}
($|b_\rho|<\delta$, $\kappa\le e^{2d\delta}=1+O(\delta)$, sample like on-line points) vs
\emph{deep} ($|b_\rho|\ge\delta$, amplified, the only danger). Thus B-2 reduces to:
deep off-line zeros do not dominate any Nyquist cell at large height.

\begin{hypothesis}[Uniform local sparsity of deep off-line zeros]\label{def:H}
There exist $\delta>0$, $\eta>0$, $T_0$ such that for all $T\ge T_0$ and every interval
$I$ of length $1/d$ centred at height $T$,
\[
\#\{\rho:\ |b_\rho|<\delta,\ t_\rho\in I\}\ \ge\ \eta\,\rho(T)\,|I| .
\]
Equivalently, deep off-line zeros occupy at most $(1-\eta)$ of each Nyquist cell.
\end{hypothesis}

\begin{lemma}[(H) $\Rightarrow$ coercivity]\label{lem:E}
Under \textnormal{(H)}, every Nyquist cell contains $\ge\eta\rho|I|$ samples with
$\kappa=1+O(\delta)$; this set has no Nyquist gaps, so by Beurling
$\sum_{\text{shallow}\cup\text{on}}|F(t_\rho)|^{2}\ge c\eta\rho\Norm{F}_2^{2}$. Taking
$\delta$ small gives $\tplus(F)\ge c'\rho\Norm{F}_2^{2}$.
\end{lemma}

\subsection{The honest status of (H)}\label{sec:Hstatus}
(H) is strictly weaker than RH: RH says there are \emph{no} off-line zeros; (H) says only
that deep ones do not \emph{locally dominate} (they may exist, even infinitely many,
provided they do not fill a Nyquist cell). It is implied by RH and by any uniform
short-interval bound $N(\half+\delta,T+\tfrac1d)-N(\half+\delta,T)=o(\rho)$.

However, (H) is \emph{not} established unconditionally by current methods. The global
density $N(\half+\delta,T)\ll T^{1-c\delta}$ leaves a budget allowing $\sim\log T$ deep
zeros to be concentrated in a single short interval (cost $\log T\ll T^{1-c\delta}$),
which the global bound does not forbid. Selberg's fluctuation theory makes such a
cell-filling cluster a large deviation of $S$, of measure
$\ll\exp(-c(\log T)^{2}/\log\log T)$---extraordinarily rare, but not provably absent at
sporadic large heights. And (H) is necessary: without it there exist $F$ with
$\tplus(F)=0$ but $\tminus(F)\ne0$, breaking closability in $\Hplus$. Alternative
realizations (all-zeros norm, $L^{2}$) reduce to the same gap/clustering condition.
Hence \textbf{Theorem~\ref{thm:A} is unconditional modulo (H), with (H) a uniform
short-interval zero-density bound far weaker than RH, very likely true, but at the
frontier of the unconditional.}

% ════════════════════════════════════════════════════════════════════════
\section{Theorem A: the semibounded realization}\label{sec:thmA}

\begin{proof}[Proof of Theorem~\ref{thm:A}]
By Lemmas~\ref{lem:C}--\ref{lem:D} and Lemma~\ref{lem:E} (under (H)),
$\ttilde\le \tfrac{C_d}{c'}\tplus=:C\tplus$, hence $|\tminus|\le 4C\tplus$ (the factor
$4$ is removed in \S\ref{sec:angular}). In $\Hplus$, where $\tplus=\Norm{\cdot}_{\Hplus}^2$,
\[
\frac{\tform(F)}{\Norm{F}_{\Hplus}^{2}}=1+\frac{\tminus(F)}{\tplus(F)}\ge 1-4C .
\]
$\tplus$ is closed (a monotone limit of bounded forms; Kato \cite{Kato}), and $\tminus$
is form-bounded relative to it with bound $4C$, so $\tform=\tplus+\tminus$ is closed and
semibounded; $\Top$ is the self-adjoint operator associated to it via KLMN, with form
domain $\Hplus$.
\end{proof}

\begin{corollary}[Faithful reformulation]\label{cor:Aprime}
The positivity $\tform\succeq0$ does not refer to the inner product, so its truth is
norm-independent; with Theorem~\ref{thm:A},
\[
\boxed{\ \mathrm{RH}\iff\inf\spec(\Top)\ge0\ }
\]
---RH is exactly the sign of the finite bottom of the rigorously realized operator $\Top$.
\end{corollary}

The magnitude estimates give $C\ge e^{d}>1$, so the bottom $1-4C$ is negative; lowering
$C$ to $\le\tfrac14$ is exactly RH. Theorem~\ref{thm:A} relocates RH to the sign of a
concrete finite object; it does not evaluate that sign.

% ════════════════════════════════════════════════════════════════════════
\section{Theorem B (i): the Kre\u{\i}n identity}\label{sec:krein}

This section is unconditional. Define the evaluation operator
$E:\PW\to\ell^{2}(\rho)$, $(EF)_\rho:=F(t_\rho+ib_\rho)$ (bounded by Lemma~\ref{lem:D}),
and the fundamental symmetry $J$ as the linear involution induced by the partner map
\eqref{eq:sigma}, $(Jc)_\rho:=c_{\sigma(\rho)}$. Then $J^{2}=I$ and $J=J^{*}$ (a
permutation by an involution is self-adjoint): a genuine Kre\u{\i}n fundamental symmetry.

\begin{proposition}[The identity]\label{prop:identity}
With $\langle a,b\rangle=\sum_\rho a_\rho\overline{b_\rho}$,
\[
\langle EF,JEF\rangle=\sum_\rho (EF)_\rho\overline{(EF)_{\sigma(\rho)}}
=\sum_\rho F(t_\rho+ib_\rho)\overline{F(t_\rho-ib_\rho)}=\tform(F),
\]
i.e. $\boxed{\ \tform=E^{*}JE,\quad J=J^{*},\ J^{2}=I\ }$ (exact, unconditional).
\end{proposition}
\begin{proof}
Immediate from $b_{\sigma(\rho)}=-b_\rho$, $t_{\sigma(\rho)}=t_\rho$ and \eqref{eq:EF}.
\end{proof}

\begin{remark}
Writing $J$ with a complex conjugation, $(Jc)_\rho=\overline{c_{\bar\rho}}$, makes it
antilinear; that coincides with the linear $J$ above on real-type test functions
$F(\bar z)=\overline{F(z)}$. The linear involution is the convention-free statement.
\end{remark}

\paragraph{Three structural consequences.}
\begin{enumerate}[label=\textbf{(C\arabic*)}]
\item \emph{RH $\Rightarrow J=I\Rightarrow\Top=E^{*}E\ge0$.} Under RH all $b_\rho=0$, so
$\sigma(\rho)=\rho$, $J=I$, $\tform=E^{*}E\ge0$: Weil positivity becomes the trivial
positivity of a Gram operator.
\item \emph{RH fails $\Rightarrow$
$J=I_{\mathrm{on}}\oplus\bigoplus_{\text{pairs}}\left(\begin{smallmatrix}0&1\\1&0\end{smallmatrix}\right)$.}
Each off-line quartet splits under $\sigma$ into two swap-pairs $\{\rho,1-\bar\rho\}$,
each a $2\times2$ swap with spectrum $\{+1,-1\}$. All negative directions of $J$ come
exactly from off-line zeros.
\item \emph{RH $\iff E^{*}JE\ge0\iff R(E)$ is a $J$-non-negative subspace.}
\end{enumerate}
The conceptual gain is the clean separation of the geometry of $\PW$ (through $R(E)$),
the location of zeros (through $\sigma$, hence $J$), and the sign (through the indefinite
form $[\cdot,\cdot]_J$).

% ════════════════════════════════════════════════════════════════════════
\section{Theorem B (ii): the angular operator and the single axis}\label{sec:angular}

Let $\ell^{2}(\rho)=\mathcal K_+\oplus\mathcal K_-$ with $J=P_+-P_-$. On a swap-pair the
$\pm1$ eigenvectors are $\tfrac1{\sqrt2}(e_\rho\pm e_{\sigma\rho})$, so with
$F_\pm:=F(t_\rho\pm ib_\rho)$,
\[
\Norm{P_+EF}^{2}=\tplus+\!\!\sum_{\text{pairs}}\tfrac12|F_++F_-|^{2},\qquad
\Norm{P_-EF}^{2}=\!\!\sum_{\text{pairs}}\tfrac12|F_+-F_-|^{2},
\]
and exactly (no factor-$4$ slack)
\begin{equation}\label{eq:split}
\boxed{\ \tform(F)=\Norm{P_+EF}^{2}-\Norm{P_-EF}^{2}.\ }
\end{equation}

\begin{lemma}[Qualitative injectivity; unconditional, not circular]\label{lem:inj}
$P_+EF=0\Rightarrow F=0$.
\end{lemma}
\begin{proof}
$P_+EF=0$ forces $F$ to vanish at all on-line zeros. By the type-density uniqueness
theorem (a nonzero $L^{2}$ function of exponential type $d$ cannot vanish on a real
sequence of density $>d/\pi$ \cite{Levin}), and the on-line zeros have density
$\tfrac1{2\pi}\log T\to\infty\gg d/\pi$, we get $F=0$. This uses only the zeros of $F$,
no sampling/coercivity input.
\end{proof}

Hence $\overline{R(E)}$ is the graph of the \emph{angular operator}
$\Kop:=P_-(P_+|_{R(E)})^{-1}:\mathcal K_+\supseteq\mathrm{dom}\,\Kop\to\mathcal K_-$,
defined unconditionally (possibly unbounded). Kre\u{\i}n theory
\cite{AzizovIokhvidov,Bognar} gives the single-axis dictionary
\begin{equation}\label{eq:dict}
\boxed{\
\underbrace{\text{type-density}}_{\Kop\ \text{defined}}\,;\quad
\underbrace{(H)\Rightarrow\Norm{\Kop}<\infty}_{\text{B-2 (finite bottom)}}\,;\quad
\underbrace{\Norm{\Kop}\le1\iff\mathrm{RH}}_{\text{given }\Kop\text{ defined}}.\ }
\end{equation}
Here $\Norm{\Kop}<\infty\iff R(E)$ uniformly positive-projecting $\iff$ the coercivity
B-2 needs; (H) is sufficient for it (the honest direction is
$(H)\Rightarrow\Norm{\Kop}<\infty\Rightarrow$ B-2). And $\Norm{\Kop}\le1\iff
\Norm{P_-EF}\le\Norm{P_+EF}\ \forall F\iff\tform\ge0\iff$ RH. From Lemma~\ref{lem:C},
$\Norm{\Kop}\ge e^{d/2}>1$ on the magnitude side---the wrong side of $1$.

\paragraph{Cleanest statement of this line of work.} B-2 is ``$\Kop$ bounded'' (needs (H)); RH
is ``$\Kop$ a contraction.''

\subsection{Where the negativity lives, and a warning}
On a swap-pair the negative coordinate is the vertical antisymmetric difference
\[
d_\rho:=F(t_\rho+ib_\rho)-F(t_\rho-ib_\rho)=2i\,b_\rho\,F'(t_\rho)+O(b_\rho^{3}\Norm{F''}),
\]
so to leading order $\Norm{P_-EF}^{2}\approx\sum_{\text{pairs}}2\,b_\rho^{2}|F'(t_\rho)|^{2}$:
the negative directions are depth-weighted derivative samples, and $\Norm{\Kop}\to0$ as
depths $\to0$.

\begin{remark}[The optimistic-pillar trap]
The $b_\rho^{2}$ form is the \emph{shallow linearization only}. The Taylor remainder is
$O(b_\rho^{3}\Norm{F''})$ with $\Norm{F''}\le d^{2}\Norm{F}$ (Bernstein); for \emph{deep}
zeros ($|b_\rho|$ up to $\half$, $d$ moderately $>\half$) the remainder/leading ratio is
$\sim b_\rho^{2}d=O(1)$---the expansion does not converge, and the true vertical
evaluation carries the $e^{2d|b_\rho|}$ amplification responsible for
$\Norm{\Kop}\ge e^{d/2}$. Thus a ``$\sum b_\rho^{2}|F'|^{2}\le\theta\tplus\Rightarrow$ RH''
inequality would deliver RH only on the shallow part already known harmless. The correct
concrete target keeps the full vertical difference:
\begin{equation}\label{eq:rayleigh}
\Norm{\Kop}^{2}=\sup_F\frac{\sum_{\text{pairs}}\tfrac12|F(t_\rho+ib_\rho)-F(t_\rho-ib_\rho)|^{2}}
{\tplus+\sum_{\text{pairs}}\tfrac12|F(t_\rho+ib_\rho)+F(t_\rho-ib_\rho)|^{2}},
\qquad \Norm{\Kop}\le1\iff\mathrm{RH}.
\end{equation}
\end{remark}

% ════════════════════════════════════════════════════════════════════════
\section{Theorem C: a no-go theorem for the zero-side attack}\label{sec:nogo}

\begin{proof}[Proof of Theorem~\ref{thm:C}]
The pivot is that the explicit formula gives $\tform(F)$ as the \emph{same} functional in
two ways:
\begin{equation}\label{eq:pivot}
\underbrace{\Norm{P_+EF}^{2}-\Norm{P_-EF}^{2}}_{\text{zero side: depends on }\{t_\rho,b_\rho\}}
=\tform(F)=
\underbrace{\rmar(F)+\rmq(F)-\rmp(F)}_{\text{arithmetic side: contains no }b_\rho}.
\end{equation}

\textbf{(i) No independent information in the depths.} The arithmetic side sums the
depths out---it does not contain $b_\rho$. Since both sides equal the same functional of
$F$, any bound on $\Norm{\Kop}$ phrased through the $b_\rho$ is, by \eqref{eq:pivot}, a
bound on $\tform$ through the arithmetic side in disguise.

\textbf{(ii) Tautology w.r.t. Weil.} $\Norm{\Kop}\le1\iff\Norm{P_-EF}\le\Norm{P_+EF}\
\forall F\iff\tform\ge0\iff\rmar+\rmq-\rmp\ge0\iff$ RH. Lowering $\Norm{\Kop}$ to $\le1$
\emph{is} proving arithmetic positivity.

\textbf{(iii) Sign-agnostic bounds land on the wrong side.} Dropping the phase
cancellation between $F_+$ and $F_-$ (triangle inequality
$|F_+-F_-|^{2}\le2(|F_+|^{2}+|F_-|^{2})$) gives
$\Norm{P_-EF}^{2}\le\kappa\cdot(\text{off-line cell count})\Norm{F}^{2}$, $\kappa$ up to
$e^{2d|b|}$, hence $\Norm{\Kop}^{2}\gtrsim e^{d}$. Retaining the phase requires the exact
analytic continuation of $F$ at $t_\rho\pm ib_\rho$; for deep zeros that is not the
$2ib_\rho F'$ linearization but the global behaviour of $F$---the arithmetic side again.
\end{proof}

\begin{corollary}[Why the zero-side realization is circular as a tool]\label{cor:circular}
To even write $\Kop=P_-(P_+|_{R(E)})^{-1}$ one must know the off-line zero data
$\{(t_\rho,b_\rho):b_\rho\ne0\}$---the very data whose emptiness is RH. The realization is
a true identity but presupposes the unknowns. The Connes--Consani realization is the
useful one precisely because its operator is built from explicit archimedean data,
independent of the zeros. This is the structural reason the zero-side picture cannot
attack the sign.
\end{corollary}

\[
\boxed{\ \begin{array}{c}
\textbf{No-go: } \Norm{\Kop}\text{ depends on the off-line depths only through }\tform;\\[2pt]
\text{bounding }\Norm{\Kop}\le1\text{ from the zero side}\equiv\text{Weil positivity}\equiv\text{RH.}
\end{array}\ }
\]

% ════════════════════════════════════════════════════════════════════════
\section{Relation to prior work}\label{sec:prior}

\paragraph{Connes--Consani \cite{ConnesConsani}.} They root Weil positivity in the
\emph{compression of the scaling action} $\vartheta$ on $L^{2}(\mathbb R)_{\mathrm{ev}}$
onto Sonin's space $\mathbf S$ (the orthogonal complement of the range of the phase-space
cutoff projection), proving $W_\infty(g*g^{*})\ge\tr(\vartheta(g)\,\mathbf S\,\vartheta(g)^{*})$
and controlling the residual difference (between the Weil distribution and the Sonin
trace) via prolate spheroidal wave functions and Hermitian Toeplitz matrices, by an
eigenvalue-below-$1$ argument. Our reduction ``$\Norm{\Kop}\le1\iff$ RH'' is of the same
type (a compression with a norm/eigenvalue-below-$1$ positivity criterion). We do
\emph{not} claim an identification $U\,\Kop\,U^{-1}=K_{\mathrm{CC}}$: the two are distinct
realizations of the same form $\tform$---ours \emph{zero-side} (depends on the zeros),
theirs \emph{archimedean} (explicit, independent of the zeros). Theorem~\ref{thm:C} and
Corollary~\ref{cor:circular} explain the asymmetry: the archimedean realization is a
genuine handle, the zero-side one is circular.

\paragraph{Suzuki \cite{Suzuki}.} Completing the Weil hermitian form yields, \emph{under
RH}, a de~Branges space $\mathcal H(E_\xi)$ with $E_\xi(z)=\xi(\half-iz)+\xi'(\half-iz)$
(Hermite--Biehler under RH \cite{Lagarias}), turning Weil's inequality into equalities.
Suzuki notes (after Selberg) that ``the construction of a space assuming RH will not be
useful for attacking RH'' and therefore builds unconditional spaces $\mathcal H_0,
\mathcal K_0$ via screw lines. Our completion is the same in spirit but \emph{directly}
unconditional: with off-line zeros allowed, the metric is genuinely indefinite, so the
natural object is a \emph{Kre\u{\i}n} space (Pontryagin only if the off-line set were
finite, which is not known). The Kre\u{\i}n-vs-de~Branges dichotomy is precisely the
positivity-built-in vs not distinction; our Theorem~\ref{thm:C} is the rigorous form of
Selberg's warning, specialized to the zero-side realization.

\paragraph{Classical lineage.} Weil's criterion \cite{Weil1952,Bombieri}; de~Branges'
\emph{Hilbert Spaces of Entire Functions} \cite{deBranges} (the positive regime);
Azizov--Iokhvidov \cite{AzizovIokhvidov} and Bogn\'ar \cite{Bognar} for Kre\u{\i}n-space
angular operators; Montgomery \cite{Montgomery} and Selberg \cite{Selberg} for the zero
statistics behind (H). Our place is a modular reformulation plus a no-go, not a new
mechanism for positivity.

% ════════════════════════════════════════════════════════════════════════
\section{Honest status and open problems}\label{sec:open}

\begin{center}
\small
\begin{tabular}{ll}
\toprule
Statement & Status \\
\midrule
Lemmas~\ref{lem:A}--\ref{lem:D} (IBP, quadrature $O(\tplus)$, vertical, P--P) & unconditional \\
Depth split; (H) $\Rightarrow$ coercivity (Lemma~\ref{lem:E}) & unconditional \\
Theorem~\ref{thm:A} (semibounded $\Top$, $\inf\spec\ge1-4C$) & \textbf{modulo (H)} \\
Corollary~\ref{cor:Aprime} (RH $\iff$ sign of the bottom) & modulo (H) \\
Theorem~\ref{thm:B} (i): $\tform=E^{*}JE$ & unconditional \\
Theorem~\ref{thm:B} (ii): angular operator, dictionary \eqref{eq:dict} & unconditional \\
Theorem~\ref{thm:C}: no-go for the zero-side attack & unconditional \\
RH itself (the sign) & not proved; Thm~\ref{thm:C} shows this route cannot \\
\bottomrule
\end{tabular}
\end{center}

\paragraph{What is not claimed.} No proof of RH. Theorem~\ref{thm:A} is conditional on
(H); (H) is much weaker than RH but not unconditionally established (\S\ref{sec:Hstatus}).
The Kre\u{\i}n picture is the correct language and the precise target, and
Theorem~\ref{thm:C} proves it is not a route past the sign.

\paragraph{Open problems.}
\begin{enumerate}[label=\textbf{\arabic*.}]
\item \emph{Is (H) unconditional?} Is the uniform (every cell, all large $T$)
short-interval zero-density for fixed $\sigma>\half$ unconditional, or only on average?
This is the one item separating ``unconditional B-2'' from ``B-2 modulo (H)''.
\item \emph{Sharp constants.} The Plancherel--P\'olya constant of Lemma~\ref{lem:D} with
clustering and margin $d-\half$; the quantitative Beurling step in Lemma~\ref{lem:E}.
\item \emph{The only route to the sign.} By Theorem~\ref{thm:C} the sign cannot come from
the zero side; it requires \emph{structural positivity}---a factorization $\Top=A^{*}A$,
Connes' archimedean trace positivity, or a de~Branges chain---an independently-given
positive structure, not a bound. This is the genuine RH problem and is left open. It is
the subject of the companion research plan (route B).
\end{enumerate}

\paragraph{Conclusion.} Starting from Weil's criterion (RH $\iff$ a positivity), we built
a faithful, RH-independent, semibounded operator realization $\Top$ whose bottom is
finite and whose sign is RH (Theorem~\ref{thm:A}), exposed its exact Kre\u{\i}n structure
$\tform=E^{*}JE$ and the single controlling axis $\Norm{\Kop}$ (Theorem~\ref{thm:B}), and
proved that this zero-side structure---however explicit its dependence on the off-line
zeros---cannot decide the sign, because the depths enter only through the very functional
whose positivity is RH (Theorem~\ref{thm:C}). The program therefore delivers a precise
localization of RH's difficulty (the contraction $\Norm{\Kop}\le1$) and a rigorous
account of why the natural zero-side attacks fail---but not a proof. The sign, as long
suspected, requires structural positivity, which the magnitude/zero-side machinery
provably cannot supply. This closes the explicit-formula/zero-side line of this line of work.

% ════════════════════════════════════════════════════════════════════════
\begin{thebibliography}{99}

\bibitem{AzizovIokhvidov} T.~Ya.~Azizov, I.~S.~Iokhvidov, \emph{Linear Operators in
Spaces with an Indefinite Metric}, Wiley, Chichester, 1989.

\bibitem{Bognar} J.~Bogn\'ar, \emph{Indefinite Inner Product Spaces}, Springer, 1974.

\bibitem{Bombieri} E.~Bombieri, \emph{Problems of the Millennium: The Riemann Hypothesis},
Clay Mathematics Institute, 2000.

\bibitem{ConnesConsani} A.~Connes, C.~Consani, \emph{Weil positivity and trace formula,
the archimedean place}, Selecta Math.\ (2021); arXiv:2006.13771.

\bibitem{deBranges} L.~de~Branges, \emph{Hilbert Spaces of Entire Functions},
Prentice-Hall, Englewood Cliffs, 1968.

\bibitem{Kato} T.~Kato, \emph{Perturbation Theory for Linear Operators}, 2nd ed.,
Springer, 1976.

\bibitem{Lagarias} J.~C.~Lagarias, \emph{Hilbert spaces of entire functions and Dirichlet
$L$-functions}, in \emph{Frontiers in Number Theory, Physics, and Geometry I}, Springer,
2006, 365--377.

\bibitem{Levin} B.~Ya.~Levin, \emph{Lectures on Entire Functions}, Transl.\ Math.\
Monographs 150, Amer.\ Math.\ Soc., 1996.

\bibitem{Montgomery} H.~L.~Montgomery, \emph{The pair correlation of zeros of the zeta
function}, in \emph{Analytic Number Theory}, Proc.\ Sympos.\ Pure Math.\ 24 (1973),
181--193.

\bibitem{Selberg} A.~Selberg, \emph{On the remainder in the formula for $N(T)$, the number
of zeros of $\zeta(s)$ in the strip $0<t<T$}, Avh.\ Norske Vid.\ Akad.\ Oslo (1944).

\bibitem{Suzuki} M.~Suzuki, \emph{On the Hilbert space derived from the Weil
distribution}, arXiv:2301.00421 (v3, 2025).

\bibitem{Weil1952} A.~Weil, \emph{Sur les ``formules explicites'' de la th\'eorie des
nombres premiers}, Comm.\ S\'em.\ Math.\ Univ.\ Lund (1952), 252--265.

\end{thebibliography}

\end{document}
